% =========================================================
%  sesi/06/sesi.tex
%  Sesi 6 — Level Lorong dan Penjaga yang Mengejar
% =========================================================

\sesi{06}{Level Lorong dan Penjaga yang Mengejar}{2.1}{130 menit}

Sampai sesi 5, seluruh permainan berlangsung pada satu garis mendatar. Pemain
berlari ke kiri dan ke kanan, gravitasi menariknya ke bawah, dan setiap rintangan
berupa rak atau peti. Pada minggu 6 ruang permainan berubah: Anda membuka level
kedua, yaitu Lorong, yang dilihat dari atas dan disusun dari petak-petak.

Perubahan pandangan itu membawa akibat yang besar. Tidak ada lagi gravitasi dan
tidak ada lagi lompatan; yang ada adalah gerak ke empat arah pada bidang datar.
Pemain berjalan dengan WASD, dinding menghalangi langkahnya, dan seekor penjaga
mengikutinya dari sudut ruangan. Penjaga itu belum tahu apa pun tentang peta.
Ia hanya tahu di mana pemain berada, lalu mengejarnya. Perilaku itu disebut
\istilah{chase} dan tampak seperti makhluk berkehendak.

Sesi ini menuntaskan Sub-CPMK 2.1: merancang perilaku NPC sederhana berdasarkan
posisi pemain. Seluruh isi sesi bermuara pada satu perhitungan vektor yang
panjangnya hanya beberapa baris, tetapi dipakai berulang pada setiap bingkai.

% ---------------------------------------------------------
\bagian{Tujuan Sesi}

\begin{tujuanbox}
Setelah menuntaskan sesi ini, Anda dapat menunjukkan:
\begin{enumerate}[leftmargin=1.5em, itemsep=0.2em]
  \item level Lorong berukuran $24 \times 14$ petak dengan sel 36 piksel yang
        tampil di jendela, lengkap dengan dinding, titik awal, dan pintu keluar;
  \item pemain yang bergerak ke empat arah dengan WASD atau tombol panah dan
        tidak dapat menembus dinding;
  \item penjaga yang mulai bergerak ketika pemain masuk radius 7 petak
        (252 piksel) dan berhenti pada jarak 26 piksel;
  \item perhitungan vektor arah, jarak, dan kecepatan penjaga yang Anda
        kerjakan di kertas dengan angka yang benar;
  \item penjelasan mengapa vektor arah harus dinormalkan sebelum dikalikan
        kecepatan, beserta satu kelemahan AI reaktif yang Anda amati sendiri.
\end{enumerate}
\end{tujuanbox}

% ---------------------------------------------------------
\bagian{Prasyarat}

\begin{prasyaratbox}
\begin{itemize}[leftmargin=1.4em, itemsep=0.25em]
  \item Checkpoint sesi 5 berjalan: pemain, artefak, peti, dan pintu sudah
        tersusun sebagai \istilah{sprite} di dalam \istilah{group}.
  \item Anda dapat memakai \code{pygame.Vector2} untuk penjumlahan,
        pengurangan, dan \code{normalize()}, sebagaimana dilatih pada sesi 3.
  \item Anda memahami pemakaian \code{delta_seconds} pada pembaruan posisi.
  \item Bab VIII buku ajar sudah dibaca, khususnya bagian tentang persepsi,
        penormalan, dan jarak henti.
\end{itemize}
\end{prasyaratbox}

% ---------------------------------------------------------
\bagian{Peta Kemampuan}

\petasesi
  {Level Rak bertampilan samping; pemain dan benda bergerak pada satu sumbu
   mendatar.}
  {Level Lorong berbasis petak, gerak empat arah dengan tabrakan dinding,
   penjaga reaktif dengan radius pandang dan jarak henti, serta garis arah.}
  {Mesin keadaan penjaga pada sesi 7 dan pencarian jalur A* pada sesi 8.}

% ---------------------------------------------------------
\bagian{Konsep yang Dipakai}

\subsection*{Petak sebagai ruang}

Peta Lorong disimpan sebagai daftar baris teks. Setiap karakter menempati satu
petak: \texttt{\#} adalah dinding, titik adalah lantai, \code{S} adalah titik awal
pemain, dan \code{E} adalah pintu keluar. Peta yang dipakai berukuran 24 kolom
dan 14 baris, dengan setiap petak berukuran 36 piksel. Pojok kiri atas peta
diletakkan pada $(48, 76)$ piksel di dalam jendela, sehingga seluruh peta
menempati $864 \times 504$ piksel dan masih menyisakan ruang untuk HUD.

Dari data teks itu program membangun dua hal. Pertama, kumpulan koordinat
petak dinding. Kedua, daftar \code{pygame.Rect} untuk keperluan tabrakan. Titik
pusat sebuah petak $(k, b)$ dihitung dengan
\[
x = 48 + 36k + 18, \qquad y = 76 + 36b + 18 .
\]
Pemisahan antara data peta dan gambar sangat membantu: mengubah lorong berarti
mengubah teks peta, bukan menyunting puluhan koordinat.

\subsection*{Pemain bergerak pada dua sumbu}

Masukan WASD menghasilkan vektor arah yang komponennya $-1$, $0$, atau $1$.
Apabila pemain menekan dua tombol sekaligus, panjang vektor itu menjadi
$\sqrt{2}$; karena itu vektor dinormalkan lebih dahulu, lalu dikalikan kecepatan
pemain 210 piksel per detik. Tanpa penormalan, gerak diagonal akan lebih cepat
daripada gerak lurus.

Posisi diperbarui satu sumbu demi satu sumbu. Sumbu $x$ digerakkan, kotak
tabrakan diselaraskan, lalu dinding diselesaikan; sesudah itu sumbu $y$
dikerjakan dengan cara yang sama. Urutan itu membuat pemain dapat menyusuri
dinding, tetapi tidak dapat menembusnya dari sudut.

\subsection*{Penjaga: dari selisih posisi ke kecepatan}

Penjaga tidak menerima arah dari luar. Setiap bingkai ia menghitung sendiri
arah menuju pemain. Misalkan $\vec{n}$ adalah posisi penjaga dan $\vec{p}$
posisi pemain, keduanya dalam piksel. Empat besaran berikut dihitung berurutan.
\begin{align*}
  \vec{u} &= \vec{p} - \vec{n}, &
  d &= \lVert \vec{u} \rVert, \\
  \hat{u} &= \frac{\vec{u}}{\lVert \vec{u} \rVert}, &
  \vec{v} &= \hat{u}\,s .
\end{align*}
Di sini $\vec{u}$ adalah vektor dari penjaga menuju pemain, $d$ adalah jaraknya
dalam piksel, $\hat{u}$ adalah vektor satuan yang panjangnya satu, dan $s$
adalah kecepatan penjaga, yaitu 165 piksel per detik. Urutan pengurangan penting:
$\vec{p} - \vec{n}$ menunjuk ke pemain, sedangkan urutan sebaliknya membuat
penjaga menjauh.

Penjaga hanya bergerak apabila $26 < d \le 252$. Angka 252 berasal dari radius
7 petak dikalikan 36 piksel, sedangkan 26 piksel adalah jarak henti. Di luar
radius itu pemain dianggap tidak terlihat; di dalam jarak henti penjaga berhenti
agar tidak bergetar di atas pemain. Pada 60 bingkai per detik, perpindahan satu
bingkai hanya 2,75 piksel, jauh lebih kecil daripada jarak henti.

Satu garis merah digambar dari penjaga ke pemain selama penjaga mengejar. Garis
itu bukan bagian permainan; ia alat periksa yang memperlihatkan arah hasil
perhitungan. Warna badan penjaga juga berubah: abu-abu ketika berdiam dan merah
ketika mengejar.

\subsection*{Reaktif dan berbasis perencanaan}

Penjaga pada sesi ini disebut reaktif karena keputusannya hanya memakai keadaan
yang terlihat pada bingkai itu, yaitu posisi pemain. Ia tidak menyimpan peta
dinding dan tidak merencanakan rute. Akibatnya, ketika pemain bersembunyi di
balik dinding, penjaga tetap menekan dinding itu selama pemain berada di
seberangnya. AI berbasis perencanaan menambahkan keterangan yang hilang itu:
ia mengetahui petak mana yang dapat dilalui, menyusun rangkaian langkah menuju
pemain, lalu memilih rute terpendek. Pencarian rute semacam itu, dengan
algoritma A*, menjadi pokok sesi 8.

% ---------------------------------------------------------
\bagian{Tutorial}

Kerjakan langkah berikut berurutan. Setiap langkah adalah satu tindakan yang
selesai; jalankan program setiap kali satu bagian baru selesai ditulis.

\begin{langkah}
  \item \textbf{Siapkan proyek.} Buka folder proyek sesi 5, lalu jalankan
        \code{python main.py} untuk memastikan level Rak masih berjalan.
        Salin folder itu menjadi folder kerja sesi 6 agar checkpoint lama tetap
        utuh.

  \item \textbf{Tambahkan pengaturan Lorong.} Pada \file{settings.py}, tambahkan
        \code{CELL = 36}, \code{GRID_OFFSET_X = 48}, \code{GRID_OFFSET_Y = 76},
        dan \code{TOP_DOWN_SPEED = 210.0}. Tambahkan pula pengaturan penjaga:
        \code{GUARD_SIZE = 30}, \code{GUARD_RADIUS_CELLS = 7.0},
        \code{GUARD_STOP_DISTANCE = 26.0}, dan \code{GUARD_CHASE_SPEED = 165.0}.

  \item \textbf{Tulis peta Lorong.} Pada \file{levels.py}, buat konstanta
        \code{SCENARIOS} berisi satu peta 24 kolom dan 14 baris. Gunakan
        \texttt{\#} untuk dinding, titik untuk lantai, \code{S} untuk titik awal,
        dan \code{E} untuk pintu keluar. Periksa bahwa setiap baris memiliki
        panjang yang sama; satu karakter yang kurang akan menggeser seluruh
        peta.

  \item \textbf{Bangun level dari teks.} Tulis fungsi \code{build_lorong} yang
        menyisir peta, mengumpulkan petak dinding ke dalam himpunan, mencari
        titik awal dan pintu keluar, lalu mengubah setiap dinding menjadi
        \code{pygame.Rect}. Kembalikan seluruhnya sebagai satu kamus.

  \item \textbf{Tambahkan titik pusat petak.} Tulis fungsi \code{cell_center}
        pada \file{entities.py}. Fungsi itu menerima pasangan $(k, b)$ dan
        mengembalikan \code{pygame.Vector2} titik pusat petak tersebut. Fungsi
        ini dipakai untuk menempatkan pemain, penjaga, dan pintu keluar.

  \item \textbf{Tulis gerak pemain tampak atas.} Tambahkan metode
        \code{update_top_down} pada kelas \code{Player}, beserta pembantu
        \code{_resolve_top_down}. Kerjakan sumbu $x$ lebih dahulu, baru
        sumbu $y$, dan bulatkan posisi ke bilangan bulat setiap kali
        \code{rect} diselaraskan.

  \item \textbf{Tulis kelas penjaga.} Tambahkan kelas \code{Guard} yang
        menyimpan posisi, menghitung vektor menuju pemain, menormalkannya, dan
        memindahkan badan penjaga. Sediakan pula metode \code{reset} untuk
        mengembalikan penjaga ke titik awalnya setelah pemain tertangkap.

  \item \textbf{Sambungkan level kedua.} Pada \file{game.py}, tambahkan
        \code{_load_lorong}, \code{_update_lorong}, dan \code{_draw_lorong}.
        \code{_update_lorong} membaca WASD, memanggil gerak pemain, memanggil
        pembaruan penjaga, memeriksa tangkapan, dan memeriksa pintu keluar.

  \item \textbf{Tambahkan pilihan level.} Pada \file{main.py}, tambahkan
        argumen \code{--level} dengan nilai 1 atau 2, lalu panggil
        \code{game.start_level(args.level)} sebelum \code{run()}.

  \item \textbf{Jalankan dan amati.} Jalankan
        \code{python main.py --level 2}. Berjalanlah dengan WASD, sentuh dinding,
        lalu dekati penjaga dan menjauhlah. Ulangi beberapa kali sambil
        memperhatikan kapan penjaga mulai bergerak dan kapan ia berhenti.
\end{langkah}

\begin{ujibox}
Setelah Langkah 10, periksa lima hal berikut.
\begin{itemize}[leftmargin=1.4em, itemsep=0.2em]
  \item Peta memenuhi jendela dengan rapi: dinding berwarna biru muda, pintu
        keluar berupa petak hijau. Bila peta terpotong, periksa offset dan
        ukuran sel.
  \item Pemain bergerak sama cepat ke empat arah, dan gerak diagonal tidak
        lebih cepat. Bila diagonal terasa lebih cepat, periksa pemanggilan
        \code{normalize()}.
  \item Pemain dapat menyusuri dinding tetapi tidak dapat menembusnya, termasuk
        ketika menekan dua tombol sekaligus pada sudut.
  \item Penjaga mulai bergerak ketika jarak pemain sekitar 252 piksel dan
        berhenti ketika jarak sekitar 26 piksel. Bila ia bergetar di atas
        pemain, periksa jarak henti dan pembatasan langkah.
  \item Garis merah hanya muncul ketika penjaga mengejar, dan arahnya selalu
        menunjuk ke pemain. Bila garis tidak muncul, periksa nilai
        \code{chasing} pada bingkai itu.
\end{itemize}
\end{ujibox}

% ---------------------------------------------------------
\bagian{Kode Sesi Ini}

Potongan berikut adalah bagian pokok checkpoint sesi 6. Berkas lengkapnya,
termasuk kelas pemain dan level Rak yang tidak berubah, tersedia pada folder
\file{code/sesi-06}. Bandingkan kode Anda dengan daftar ini setelah program
berjalan, bukan sebelumnya.

\tampilkodebaris[Python]{code/sesi-06/gudang_artefak/settings.py}{41}{52}
  {Pengaturan level Lorong dan penjaga}

\tampilkodebaris[Python]{code/sesi-06/gudang_artefak/entities.py}{89}{115}
  {Gerak pemain tampak atas dan penyelesaian tabrakan per sumbu}

\tampilkodebaris[Python]{code/sesi-06/gudang_artefak/entities.py}{186}{225}
  {Kelas penjaga: persepsi, keputusan, dan gerak}

\tampilkodebaris[Python]{code/sesi-06/gudang_artefak/game.py}{137}{162}
  {Memperbarui level Lorong: masukan, penjaga, dan syarat selesai}

\tampilkodebaris[Python]{code/sesi-06/gudang_artefak/game.py}{199}{212}
  {Menggambar dinding, pintu keluar, dan garis arah}

\checkpoint{06}{Program pada folder itu dapat dijalankan dengan
\code{python main.py --level 2}; pemain bergerak dengan WASD, dinding menahan
langkahnya, dan penjaga mengejar ketika pemain berada dalam radius 252 piksel.}

% ---------------------------------------------------------
\bagian{Kesalahan Umum}

\begin{galatbox}
\begin{itemize}[leftmargin=1.4em, itemsep=0.35em]
  \item \textbf{Penjaga lambat ketika jauh dan cepat ketika dekat.}
        Gejala: kecepatan penjaga berubah mengikuti jarak. Penyebabnya, vektor
        \code{to_player} dipakai langsung sebagai kecepatan tanpa dinormalkan.
        Periksa baris \code{to_player.normalize() * S.GUARD_CHASE_SPEED}.
  \item \textbf{Penjaga bergetar di atas pemain.}
        Gejala: penjaga melompat-lompat kecil di sekitar pemain. Penyebabnya,
        jarak henti terlalu kecil atau perpindahan satu bingkai lebih besar
        daripada jarak henti. Periksa \code{GUARD_STOP_DISTANCE} dan
        pembatasan \code{if step.length() > distance}.
  \item \textbf{Pemain menembus dinding pada sudut.}
        Gejala: pemain keluar dari peta ketika menekan dua tombol sekaligus.
        Penyebabnya, kedua sumbu diperbarui sebelum tabrakan diselesaikan.
        Periksa urutan pada \code{update_top_down}: satu sumbu, selesaikan,
        lalu sumbu berikutnya.
  \item \textbf{\code{AttributeError: 'Player' object has no attribute
        'update_top_down'}.}
        Gejala: program berhenti ketika level 2 dimuat. Penyebabnya, metode baru
        ditulis di luar kelas atau namanya berbeda. Periksa indentasi dan ejaan
        nama metode pada \file{entities.py}.
  \item \textbf{\code{KeyError: 2} ketika level dimuat.}
        Gejala: peta tidak ditemukan. Penyebabnya, kunci \code{SCENARIOS} atau
        \code{GUARD_START} tidak memuat nomor skenario yang diminta. Periksa
        kedua kamus itu dan nilai bawaan \code{build_lorong}.
  \item \textbf{Garis arah tidak pernah muncul.}
        Gejala: penjaga bergerak tetapi garis merah tidak tampak. Penyebabnya,
        \code{chasing} bernilai salah karena radius dihitung dua kali, atau
        garis digambar sebelum sprite. Periksa
        \code{GUARD_RADIUS_CELLS * S.CELL} dan urutan penggambaran.
\end{itemize}
\end{galatbox}

% ---------------------------------------------------------
\clearpage
\bagian{Lembar Kerja}

Isi tabel berikut pada catatan praktikum Anda. Kolom terakhir diisi setelah
Anda membandingkan hasil pengukuran dengan nilai pada \file{settings.py}.

{\small
\renewcommand{\arraystretch}{1.4}
\begin{tabularx}{\linewidth}{@{}C{0.8cm} L{6.2cm} Y@{}}
\toprule
\textbf{No} & \textbf{Hasil yang dicatat} & \textbf{Keterangan Anda} \\
\midrule
1 & Jarak pemain ketika penjaga mulai bergerak & \\[0.35cm]
2 & Jarak pemain ketika penjaga berhenti & \\[0.35cm]
3 & Kecepatan pemain lurus dan kecepatan pemain diagonal & \\[0.35cm]
4 & Kecepatan penjaga yang Anda perkirakan pada jarak 100, 180, dan 240 piksel & \\[0.35cm]
5 & Perilaku penjaga ketika pemain bersembunyi di balik dinding & \\[0.9cm]
6 & Satu kendala yang Anda hadapi dan cara menyelesaikannya & \\[0.9cm]
\bottomrule
\end{tabularx}
}

% ---------------------------------------------------------
\bagian{Checklist Sesi}

\begin{checklistbox}
\begin{checklist}
  \item Saya dapat menunjukkan level Lorong yang utuh, lengkap dengan dinding
        dan pintu keluar.
  \item Saya dapat menggerakkan pemain ke empat arah dan menjelaskan mengapa
        arah diagonal harus dinormalkan.
  \item Saya dapat menjelaskan urutan penyelesaian tabrakan per sumbu.
  \item Saya dapat menghitung $\vec{u}$, $d$, $\hat{u}$, dan $\vec{v}$ untuk
        satu pasangan posisi penjaga dan pemain.
  \item Saya dapat menjelaskan peran radius 252 piksel dan jarak henti
        26 piksel.
  \item Saya dapat menyebutkan satu kelemahan AI reaktif yang saya amati
        sendiri pada level Lorong.
\end{checklist}
\end{checklistbox}

% ---------------------------------------------------------
\bagian{Pengait ke Sesi Berikutnya}

Penjaga pada sesi ini hanya memiliki dua kemungkinan: mengejar atau berdiam.
Keduanya disimpan pada satu variabel bernilai benar atau salah, yaitu
\code{chasing}. Pada sesi 7 variabel itu diganti mesin keadaan yang
sesungguhnya, dengan tiga \istilah{state} yang berbeda---berdiam, berpatroli,
dan mengejar---beserta \istilah{transition} yang syaratnya diperiksa satu per
satu. Penjaga akan berjalan mengelilingi titik-titik patroli, mengejar ketika
pemain terlihat, dan menyerah setelah kehilangan jejak selama waktu tertentu.
Dengan begitu, perilakunya tidak lagi hanya reaktif, melainkan memiliki
riwayat dan jadwal.
